% GENERATED FILE. Edit canonical lessons, then run npm run generate:books.
% canonical-source-hash: fnv1a64:6e4adc1e04946991
% canonical-lessons: GU-C168-masalvu, GU-C168-velan, GU-C168-vanvu, GU-C168-tavo, GU-C168-sekvu

\chapter{Kneading and rolling rotli}
\label{ch:gu-kneading-and-rolling-rotli}
\hlchaptermodality{168}

\begin{chapteropening}
\textbf{By the end of this chapter:} \emph{I can describe how rotli is kneaded, rolled and roasted.}

Complete the last lesson of chapter 168: I can describe how rotli is kneaded, rolled and roasted.
\end{chapteropening}

\section[\texorpdfstring{\gu{મસળવું} (masaḷvũ)}{મસળવું (masaḷvũ)}]{\gu{મસળવું} (masaḷvũ) — to knead, to mash}

\label{lesson:GU-C168-masalvu}



\begin{quote}
Before the new one: say the Gujarati for to steam, then the Gujarati for to fry.
\end{quote}

\subsection*{You'll want to know: \gu{મસળવું}}
\textbf{\gu{મસળવું}} — \emph{masaḷvũ} — \textquotedblleft{}to knead, to mash\textquotedblright{}.

\begin{quote}
\textbf{\gu{તે} \gu{લોટ} \gu{મસળે} \gu{છે}.} — \emph{te loṭ masaḷe che} — \textquotedblleft{}She kneads the dough.\textquotedblright{}
\end{quote}

\begin{cousinweb}
\textbf{\gu{મસળવું}} (\emph{masaḷvũ}) is a tadbhava. CDIAL lists it, glossed \textquotedblleft{}to rub, crumble, knead\textquotedblright{}, under Sanskrit \emph{maṣati}, \textquotedblleft{}rubs, crushes\textquotedblright{}, extended with \emph{-l-} (CDIAL 9919).
\end{cousinweb}

\begin{grammarlens}[title={What you've built}]
One more word for everyday life at A2.
\end{grammarlens}

\subsection*{Your turn}
\begin{itemize}
\raggedright
  \item \emph{Say it:} \emph{masaḷvũ}
  \item \emph{Say it:} \emph{masaḷvũ}, once more
  \item \emph{Say it:} \emph{te loṭ masaḷe che} — She kneads the dough.
  \item \emph{Recall:} say the Gujarati for to steam, then the Gujarati for to fry, then say \emph{masaḷvũ} again
\end{itemize}

\begin{tcolorbox}[breakable,colback=teal!4,colframe=teal!35!black,title={Before you move on}]
What does \gu{મસળવું} mean? (\textquotedblleft{}To knead, to mash\textquotedblright{}.) Say it once more.
\end{tcolorbox}
\section[\texorpdfstring{\gu{વેલણ} (velaṇ)}{વેલણ (velaṇ)}]{\gu{વેલણ} (velaṇ) — rolling pin}

\label{lesson:GU-C168-velan}



\begin{quote}
Before the new one: say the Gujarati for to fry, then the Gujarati for to knead.
\end{quote}

\subsection*{You'll want to know: \gu{વેલણ}}
\textbf{\gu{વેલણ}} — \emph{velaṇ} — \textquotedblleft{}rolling pin\textquotedblright{}.

\begin{quote}
\textbf{\gu{આ} \gu{વેલણ} \gu{નવું} \gu{છે}.} — \emph{ā velaṇ navũ che} — \textquotedblleft{}This rolling pin is new.\textquotedblright{}
\end{quote}

\begin{cousinweb}
\textbf{\gu{વેલણ}} (\emph{velaṇ}) is a tadbhava. It comes from Sanskrit \emph{vēllana}, \textquotedblleft{}shaking, tossing; rolling-pin\textquotedblright{}, and CDIAL glosses the Gujarati word \textquotedblleft{}rolling-pin\textquotedblright{} (CDIAL 12120).
\end{cousinweb}

\begin{grammarlens}[title={What you've built}]
One more word for everyday life at A2.
\end{grammarlens}

\subsection*{Your turn}
\begin{itemize}
\raggedright
  \item \emph{Say it:} \emph{velaṇ}
  \item \emph{Say it:} \emph{velaṇ}, once more
  \item \emph{Say it:} \emph{ā velaṇ navũ che} — This rolling pin is new.
  \item \emph{Recall:} say the Gujarati for to fry, then the Gujarati for to knead, then say \emph{velaṇ} again
\end{itemize}

\begin{tcolorbox}[breakable,colback=teal!4,colframe=teal!35!black,title={Before you move on}]
What does \gu{વેલણ} mean? (\textquotedblleft{}Rolling pin\textquotedblright{}.) Say it once more.
\end{tcolorbox}
\section[\texorpdfstring{\gu{વણવું} (vaṇvũ)}{વણવું (vaṇvũ)}]{\gu{વણવું} (vaṇvũ) — to roll out}

\label{lesson:GU-C168-vanvu}



\begin{quote}
Before the new one: say the Gujarati for to knead, then the Gujarati for rolling pin.
\end{quote}

\subsection*{You'll want to know: \gu{વણવું}}
\textbf{\gu{વણવું}} — \emph{vaṇvũ} — \textquotedblleft{}to roll out\textquotedblright{}.

\begin{quote}
\textbf{\gu{તે} \gu{વેલણથી} \gu{રોટલી} \gu{વણે} \gu{છે}.} — \emph{te velaṇthī roṭlī vaṇe che} — \textquotedblleft{}She rolls out the rotli with a rolling pin.\textquotedblright{}
\end{quote}

\begin{cousinweb}
The origin of \textbf{\gu{વણવું}} (\emph{vaṇvũ}) in this sense is not settled. CDIAL lists Gujarati \emph{vaṇvũ} only as \textquotedblleft{}weaves\textquotedblright{}, under \emph{*vināti}, and does not list the sense of rolling dough (CDIAL 11773).
\end{cousinweb}

\begin{grammarlens}[title={What you've built}]
One more word for everyday life at A2.
\end{grammarlens}

\subsection*{Your turn}
\begin{itemize}
\raggedright
  \item \emph{Say it:} \emph{vaṇvũ}
  \item \emph{Say it:} \emph{vaṇvũ}, once more
  \item \emph{Say it:} \emph{te velaṇthī roṭlī vaṇe che} — She rolls out the rotli with a rolling pin.
  \item \emph{Recall:} say the Gujarati for to knead, then the Gujarati for rolling pin, then say \emph{vaṇvũ} again
\end{itemize}

\begin{tcolorbox}[breakable,colback=teal!4,colframe=teal!35!black,title={Before you move on}]
What does \gu{વણવું} mean? (\textquotedblleft{}To roll out\textquotedblright{}.) Say it once more.
\end{tcolorbox}
\section[\texorpdfstring{\gu{તવો} (tavo)}{તવો (tavo)}]{\gu{તવો} (tavo) — griddle}

\label{lesson:GU-C168-tavo}



\begin{quote}
Before the new one: say the Gujarati for rolling pin, then the Gujarati for to roll out.
\end{quote}

\subsection*{You'll want to know: \gu{તવો}}
\textbf{\gu{તવો}} — \emph{tavo} — \textquotedblleft{}griddle\textquotedblright{}.

\begin{quote}
\textbf{\gu{તવો} \gu{ગરમ} \gu{છે}.} — \emph{tavo garam che} — \textquotedblleft{}The griddle is hot.\textquotedblright{}
\end{quote}

\begin{cousinweb}
\textbf{\gu{તવો}} (\emph{tavo}) is a tadbhava. It comes from \emph{*tapaka}, \textquotedblleft{}frying pan\textquotedblright{}, a reconstructed form on the root \emph{tap}, \textquotedblleft{}to heat\textquotedblright{}, and CDIAL glosses the Gujarati word \textquotedblleft{}griddle\textquotedblright{} (CDIAL 5670).
\end{cousinweb}

\begin{grammarlens}[title={What you've built}]
One more word for everyday life at A2.
\end{grammarlens}

\subsection*{Your turn}
\begin{itemize}
\raggedright
  \item \emph{Say it:} \emph{tavo}
  \item \emph{Say it:} \emph{tavo}, once more
  \item \emph{Say it:} \emph{tavo garam che} — The griddle is hot.
  \item \emph{Recall:} say the Gujarati for rolling pin, then the Gujarati for to roll out, then say \emph{tavo} again
\end{itemize}

\begin{tcolorbox}[breakable,colback=teal!4,colframe=teal!35!black,title={Before you move on}]
What does \gu{તવો} mean? (\textquotedblleft{}Griddle\textquotedblright{}.) Say it once more.
\end{tcolorbox}
\section[\texorpdfstring{\gu{શેકવું} (śekvũ)}{શેકવું (śekvũ)}]{\gu{શેકવું} (śekvũ) — to roast, to toast}

\label{lesson:GU-C168-sekvu}



\begin{quote}
Before the new one: say the Gujarati for to roll out, then the Gujarati for griddle.
\end{quote}

\subsection*{You'll want to know: \gu{શેકવું}}
\textbf{\gu{શેકવું}} — \emph{śekvũ} — \textquotedblleft{}to roast, to toast\textquotedblright{}.

\begin{quote}
\textbf{\gu{તે} \gu{ધીમે} \gu{રોટલી} \gu{શેકે} \gu{છે}.} — \emph{te dhīme roṭlī śeke che} — \textquotedblleft{}She roasts the rotli slowly.\textquotedblright{}
\end{quote}

\begin{cousinweb}
\textbf{\gu{શેકવું}} (\emph{śekvũ}) is a tadbhava. It comes from \emph{*sēkk}, \textquotedblleft{}heat, foment, roast\textquotedblright{}, a reconstructed form; CDIAL lists the Gujarati \emph{sekvũ}, \textquotedblleft{}to foment, bake\textquotedblright{} (CDIAL 13581).
\end{cousinweb}

\begin{grammarlens}[title={What you've built}]
One more word for everyday life at A2.
\end{grammarlens}

\subsection*{Your turn}
\begin{itemize}
\raggedright
  \item \emph{Say it:} \emph{śekvũ}
  \item \emph{Say it:} \emph{śekvũ}, once more
  \item \emph{Say it:} \emph{te dhīme roṭlī śeke che} — She roasts the rotli slowly.
  \item \emph{Recall:} say the Gujarati for to roll out, then the Gujarati for griddle, then say \emph{śekvũ} again
\end{itemize}

\begin{tcolorbox}[breakable,colback=teal!4,colframe=teal!35!black,title={Before you move on}]
What does \gu{શેકવું} mean? (\textquotedblleft{}To roast, to toast\textquotedblright{}.) Say it once more.
\end{tcolorbox}
